\chapter{Der Rust-Simulationsserver}\label{kap:rust}

Die Simulation lief früher in R mit C++-Erweiterungen (Rcpp). Heute
übernimmt sie ein eigener Rust-Server
(\datei{league-simulator-rust/}), den R über HTTP anspricht. Der Server ist
rund hundertmal schneller als der Vorgänger und die \emph{einzige} Stelle,
an der die Modellkonstanten stehen.

\section{Bausteine}

\begin{center}
\small
\begin{tabular}{@{}lll@{}}
\toprule
Crate & Version & Zweck \\
\midrule
\code{axum}       & 0.8.9  & HTTP-Server und Routing \\
\code{tokio}      & 1.53.1 & asynchrone Laufzeit \\
\code{rayon}      & 1.12.0 & parallele Iterationen (Work-Stealing) \\
\code{rand}       & 0.10.3 & Zufallszahlen \\
\code{statrs}     & 0.19.1 & Poisson-Verteilungsfunktion \\
\code{serde}, \code{serde\_json} & 1.0 & JSON \\
\bottomrule
\end{tabular}
\end{center}

Gebaut wird mit \code{rust:1.98-alpine} gegen musl, im Release-Profil mit
\code{opt-level=3}, Link-Time-Optimierung, einer Codegen-Einheit und
\code{panic="abort"}. Der Code gliedert sich in die Module \code{api},
\code{elo}, \code{simulation} (Spiel und Saison), \code{monte\_carlo},
\code{league\_details} und \code{models}.

\section{Endpunkte}

Der Server lauscht auf Port 8080 und nimmt Anfragen bis 2~MB an
(\datei{src/api/mod.rs:15--24}).

\begin{center}
\small
\begin{tabularx}{\textwidth}{@{}lX@{}}
\toprule
Endpunkt & Zweck \\
\midrule
\code{GET /health} & Lebenszeichen für den Docker-Healthcheck: \code{\{"status":"ok", "version":…\}} \\
\code{POST /simulate} & Monte-Carlo-Simulation einer Liga \\
\code{POST /simulate/batch} & mehrere Ligen in einem Aufruf \\
\code{POST /league-details} & deterministischer ELO-Walk und Wahrscheinlichkeiten je Spiel, ohne Zufall \\
\code{POST /match-preview} & ein virtuelles Einzelspiel aus zwei ELO-Werten, für die Aufstiegsspiele \\
\bottomrule
\end{tabularx}
\end{center}

\subsection*{\texttt{POST /simulate}}

\begin{lstlisting}
{
  "schedule":   [[0, 1, 2, 1], [2, 3, null, null], ...],  // Heim, Gast, Tore, Tore
  "elo_values": [2057.26, 1876.14, ...],                    // Start-ELO je Team
  "team_names": ["FCB", "BVB", ...],
  "iterations": 10000,
  "mod_factor": 20,
  "elo_neutral":     [false, true, ...],   // optional, parallel zu schedule
  "adj_points":      [0, -50, ...],        // optional, Malus-Lauf
  "tore_slope": ..., "tore_intercept": ..., // nur Frauen-Ligen
  "group_of_team": [...], "relegation_places": 4   // nur 3. Liga
}
\end{lstlisting}

Mannschaften sind ab 0 nummeriert, fehlende Tore (\code{null}) werden
simuliert. Die Antwort enthält die Wahrscheinlichkeitsmatrix
(Mannschaft~$\times$~Platz, sortiert nach erwartetem Platz), die
Mannschaftsnamen, die Zahl der Simulationen, die Laufzeit und für die
3.~Liga die Zählung \code{relegation\_group\_counts}. Ungültige Anfragen
beantwortet der Server mit Status 400 und einer Klartextmeldung. Mehr als
\num{100000} Iterationen lehnt er ab.

\subsection*{\texttt{POST /league-details}}

Die Anfrage gleicht der von \code{/simulate}, die Mannschaften sind hier
aber ab 1 nummeriert. \code{max\_goals} (Standard 6) bestimmt die Größe der
Ergebnis-Matrix. Die Antwort enthält je Spiel: ob es gespielt ist, die Tore,
die ELO beider Mannschaften vor dem Spiel, die ELO-Änderung, beide
Torerwartungen, $P(\text{Heimsieg})$, $P(\text{Remis})$,
$P(\text{Auswärtssieg})$ und die Ergebnis-Matrix. Dazu kommen die aktuellen
ELO-Werte aller Mannschaften. Aus dieser einen Antwort entstehen Ligatabelle,
Rückblick, Ausblick und ELO-Verlauf.

\section{Wer welche Parameter sendet}

R sendet den Heimvorteil nie, und das Tormodell nur für die Frauen-Ligen.
Fehlt ein Wert, nimmt der Server seinen Standard. Damit gibt es für jede
Konstante genau eine Quelle (ADR~0002). Ein Wächtertest
(\datei{tests/testthat/test-waechter-quelltext.R}) liest den R- und
JavaScript-Quelltext als Text und schlägt fehl, sobald dort eine der
Konstanten auftaucht. R serialisiert JSON mit \code{digits = NA}, damit
\code{tore\_slope} nicht auf vier Nachkommastellen gerundet ankommt.

\section{Leistung}

Gemessen am 7.~September 2026 auf einem Entwicklerrechner, 18 Mannschaften,
voller Spielplan ohne gespielte Partien (der teuerste Fall):

\begin{center}
\small
\begin{tabular}{@{}lr@{}}
\toprule
Messgröße & Wert \\
\midrule
\num{10000} Iterationen & ca. 42~ms \\
\num{50000} Iterationen & ca. 218~ms \\
eine Saison & ca. 4~µs \\
HTTP und JSON auf R-Seite & ca. 20~ms \\
alle zehn Ligen je Zyklus & ca. 0{,}5~s \\
\bottomrule
\end{tabular}
\end{center}

Die Dauer eines Zyklus bestimmen damit die API-Abrufe, nicht die
Simulation.
